\documentclass[12pt,a4paper]{exam}
\usepackage[utf8]{inputenc}
\usepackage{amsmath,amssymb,amsthm}
\usepackage[margin=1in]{geometry}
\usepackage{graphicx}
\usepackage[colorlinks=true]{hyperref}
\pagestyle{headandfoot}
\firstpageheader{MIT OpenCourseWare}{}{\textbf{EXTREMELY HARD -- MIT LEVEL}}
\runningheader{MIT OpenCourseWare}{}{\textbf{EXTREMELY HARD -- MIT LEVEL}}
\firstpagefooter{}{Page \thepage}{}
\runningfooter{}{Page \thepage}{}

\begin{document}

\begin{center}
{\LARGE \textbf{MIT OpenCourseWare}}\\21W.755\\Writing and Reading Short Stories
\vspace{0.5cm}
{21W.755} --- {Writing and Reading Short Stories}
\vspace{0.3cm}
May 2026
\vspace{0.3cm}
\textbf{EXTREMELY HARD -- MIT LEVEL}
\end{center}

\begin{center}
\textbf{Time Limit: 3 hours}
\end{center}

\begin{center}
\textbf{Total Marks: 100}
\end{center}

\vspace{0.5cm}

\begin{questions}

\question[5]
Analyze the limitations of standard approaches to Narrative Structure when applied to edge cases in 21W.755. Construct explicit counterexamples showing where intuitive reasoning fails.

\vspace{0.3cm}

\question[5]
Construct an original proof-level problem involving Character Development for 21W.755 (Writing and Reading Short Stories) and provide a complete solution. Your problem should be at the level of a graduate qualifying exam.

\vspace{0.3cm}

\question[3]
Prove the fundamental theorem concerning Voice in the context of 21W.755 (Writing and Reading Short Stories). State the theorem precisely, provide a complete proof, and discuss three important corollaries.

\vspace{0.3cm}

\question[2]
Analyze the limitations of standard approaches to Setting when applied to edge cases in 21W.755. Construct explicit counterexamples showing where intuitive reasoning fails.

\vspace{0.3cm}

\question[5]
Analyze the limitations of standard approaches to Dialogue when applied to edge cases in 21W.755. Construct explicit counterexamples showing where intuitive reasoning fails.

\vspace{0.3cm}

\question[3]
Analyze the limitations of standard approaches to Narrative Structure when applied to edge cases in 21W.755. Construct explicit counterexamples showing where intuitive reasoning fails.

\vspace{0.3cm}

\question[2]
Prove the fundamental theorem concerning Character Development in the context of 21W.755 (Writing and Reading Short Stories). State the theorem precisely, provide a complete proof, and discuss three important corollaries.

\vspace{0.3cm}

\question[3]
Analyze the limitations of standard approaches to Voice when applied to edge cases in 21W.755. Construct explicit counterexamples showing where intuitive reasoning fails.

\vspace{0.3cm}

\question[2]
Prove the fundamental theorem concerning Setting in the context of 21W.755 (Writing and Reading Short Stories). State the theorem precisely, provide a complete proof, and discuss three important corollaries.

\vspace{0.3cm}

\question[5]
Prove the fundamental theorem concerning Dialogue in the context of 21W.755 (Writing and Reading Short Stories). State the theorem precisely, provide a complete proof, and discuss three important corollaries.

\vspace{0.3cm}

\question[3]
Analyze the limitations of standard approaches to Narrative Structure when applied to edge cases in 21W.755. Construct explicit counterexamples showing where intuitive reasoning fails.

\vspace{0.3cm}

\question[5]
Construct an original proof-level problem involving Character Development for 21W.755 (Writing and Reading Short Stories) and provide a complete solution. Your problem should be at the level of a graduate qualifying exam.

\vspace{0.3cm}

\question[4]
Construct an original proof-level problem involving Voice for 21W.755 (Writing and Reading Short Stories) and provide a complete solution. Your problem should be at the level of a graduate qualifying exam.

\vspace{0.3cm}

\question[3]
Prove the fundamental theorem concerning Setting in the context of 21W.755 (Writing and Reading Short Stories). State the theorem precisely, provide a complete proof, and discuss three important corollaries.

\vspace{0.3cm}

\question[6]
Construct an original proof-level problem involving Dialogue for 21W.755 (Writing and Reading Short Stories) and provide a complete solution. Your problem should be at the level of a graduate qualifying exam.

\vspace{0.3cm}

\question[4]
Prove the fundamental theorem concerning Narrative Structure in the context of 21W.755 (Writing and Reading Short Stories). State the theorem precisely, provide a complete proof, and discuss three important corollaries.

\vspace{0.3cm}

\question[7]
Construct an original proof-level problem involving Character Development for 21W.755 (Writing and Reading Short Stories) and provide a complete solution. Your problem should be at the level of a graduate qualifying exam.

\vspace{0.3cm}

\question[5]
Construct an original proof-level problem involving Voice for 21W.755 (Writing and Reading Short Stories) and provide a complete solution. Your problem should be at the level of a graduate qualifying exam.

\vspace{0.3cm}

\question[3]
Prove the fundamental theorem concerning Setting in the context of 21W.755 (Writing and Reading Short Stories). State the theorem precisely, provide a complete proof, and discuss three important corollaries.

\vspace{0.3cm}

\question[4]
Analyze the limitations of standard approaches to Dialogue when applied to edge cases in 21W.755. Construct explicit counterexamples showing where intuitive reasoning fails.

\vspace{0.3cm}

\question[5]
Construct an original proof-level problem involving Narrative Structure for 21W.755 (Writing and Reading Short Stories) and provide a complete solution. Your problem should be at the level of a graduate qualifying exam.

\vspace{0.3cm}

\question[5]
Prove the fundamental theorem concerning Character Development in the context of 21W.755 (Writing and Reading Short Stories). State the theorem precisely, provide a complete proof, and discuss three important corollaries.

\vspace{0.3cm}

\question[2]
Prove the fundamental theorem concerning Voice in the context of 21W.755 (Writing and Reading Short Stories). State the theorem precisely, provide a complete proof, and discuss three important corollaries.

\vspace{0.3cm}

\question[4]
Construct an original proof-level problem involving Setting for 21W.755 (Writing and Reading Short Stories) and provide a complete solution. Your problem should be at the level of a graduate qualifying exam.

\vspace{0.3cm}

\question[5]
Construct an original proof-level problem involving Dialogue for 21W.755 (Writing and Reading Short Stories) and provide a complete solution. Your problem should be at the level of a graduate qualifying exam.

\vspace{0.3cm}

\end{questions}

\newpage
\section*{Solutions}

\textbf{Solution 1:}

The edge case analysis reveals the subtle assumptions underlying standard treatments of Narrative Structure. By constructing explicit counterexamples, we see that the usual hypotheses cannot be weakened without loss of the theorem's conclusion.

\vspace{0.5cm}

\textbf{Solution 2:}

Solution approach: First recall the definitions and key theorems related to Character Development from 21W.755. Then apply the appropriate techniques to construct a rigorous argument. The solution must be self-contained and reference only the material from the course.

\vspace{0.5cm}

\textbf{Solution 3:}

The proof follows from the definitions and properties established in 21W.755. The key insight is to apply the relevant theorem and verify the hypotheses hold. Detailed solution depends on the specific formulation of the theorem.

\vspace{0.5cm}

\textbf{Solution 4:}

The edge case analysis reveals the subtle assumptions underlying standard treatments of Setting. By constructing explicit counterexamples, we see that the usual hypotheses cannot be weakened without loss of the theorem's conclusion.

\vspace{0.5cm}

\textbf{Solution 5:}

The edge case analysis reveals the subtle assumptions underlying standard treatments of Dialogue. By constructing explicit counterexamples, we see that the usual hypotheses cannot be weakened without loss of the theorem's conclusion.

\vspace{0.5cm}

\textbf{Solution 6:}

The edge case analysis reveals the subtle assumptions underlying standard treatments of Narrative Structure. By constructing explicit counterexamples, we see that the usual hypotheses cannot be weakened without loss of the theorem's conclusion.

\vspace{0.5cm}

\textbf{Solution 7:}

The proof follows from the definitions and properties established in 21W.755. The key insight is to apply the relevant theorem and verify the hypotheses hold. Detailed solution depends on the specific formulation of the theorem.

\vspace{0.5cm}

\textbf{Solution 8:}

The edge case analysis reveals the subtle assumptions underlying standard treatments of Voice. By constructing explicit counterexamples, we see that the usual hypotheses cannot be weakened without loss of the theorem's conclusion.

\vspace{0.5cm}

\textbf{Solution 9:}

The proof follows from the definitions and properties established in 21W.755. The key insight is to apply the relevant theorem and verify the hypotheses hold. Detailed solution depends on the specific formulation of the theorem.

\vspace{0.5cm}

\textbf{Solution 10:}

The proof follows from the definitions and properties established in 21W.755. The key insight is to apply the relevant theorem and verify the hypotheses hold. Detailed solution depends on the specific formulation of the theorem.

\vspace{0.5cm}

\textbf{Solution 11:}

The edge case analysis reveals the subtle assumptions underlying standard treatments of Narrative Structure. By constructing explicit counterexamples, we see that the usual hypotheses cannot be weakened without loss of the theorem's conclusion.

\vspace{0.5cm}

\textbf{Solution 12:}

Solution approach: First recall the definitions and key theorems related to Character Development from 21W.755. Then apply the appropriate techniques to construct a rigorous argument. The solution must be self-contained and reference only the material from the course.

\vspace{0.5cm}

\textbf{Solution 13:}

Solution approach: First recall the definitions and key theorems related to Voice from 21W.755. Then apply the appropriate techniques to construct a rigorous argument. The solution must be self-contained and reference only the material from the course.

\vspace{0.5cm}

\textbf{Solution 14:}

The proof follows from the definitions and properties established in 21W.755. The key insight is to apply the relevant theorem and verify the hypotheses hold. Detailed solution depends on the specific formulation of the theorem.

\vspace{0.5cm}

\textbf{Solution 15:}

Solution approach: First recall the definitions and key theorems related to Dialogue from 21W.755. Then apply the appropriate techniques to construct a rigorous argument. The solution must be self-contained and reference only the material from the course.

\vspace{0.5cm}

\textbf{Solution 16:}

The proof follows from the definitions and properties established in 21W.755. The key insight is to apply the relevant theorem and verify the hypotheses hold. Detailed solution depends on the specific formulation of the theorem.

\vspace{0.5cm}

\textbf{Solution 17:}

Solution approach: First recall the definitions and key theorems related to Character Development from 21W.755. Then apply the appropriate techniques to construct a rigorous argument. The solution must be self-contained and reference only the material from the course.

\vspace{0.5cm}

\textbf{Solution 18:}

Solution approach: First recall the definitions and key theorems related to Voice from 21W.755. Then apply the appropriate techniques to construct a rigorous argument. The solution must be self-contained and reference only the material from the course.

\vspace{0.5cm}

\textbf{Solution 19:}

The proof follows from the definitions and properties established in 21W.755. The key insight is to apply the relevant theorem and verify the hypotheses hold. Detailed solution depends on the specific formulation of the theorem.

\vspace{0.5cm}

\textbf{Solution 20:}

The edge case analysis reveals the subtle assumptions underlying standard treatments of Dialogue. By constructing explicit counterexamples, we see that the usual hypotheses cannot be weakened without loss of the theorem's conclusion.

\vspace{0.5cm}

\textbf{Solution 21:}

Solution approach: First recall the definitions and key theorems related to Narrative Structure from 21W.755. Then apply the appropriate techniques to construct a rigorous argument. The solution must be self-contained and reference only the material from the course.

\vspace{0.5cm}

\textbf{Solution 22:}

The proof follows from the definitions and properties established in 21W.755. The key insight is to apply the relevant theorem and verify the hypotheses hold. Detailed solution depends on the specific formulation of the theorem.

\vspace{0.5cm}

\textbf{Solution 23:}

The proof follows from the definitions and properties established in 21W.755. The key insight is to apply the relevant theorem and verify the hypotheses hold. Detailed solution depends on the specific formulation of the theorem.

\vspace{0.5cm}

\textbf{Solution 24:}

Solution approach: First recall the definitions and key theorems related to Setting from 21W.755. Then apply the appropriate techniques to construct a rigorous argument. The solution must be self-contained and reference only the material from the course.

\vspace{0.5cm}

\textbf{Solution 25:}

Solution approach: First recall the definitions and key theorems related to Dialogue from 21W.755. Then apply the appropriate techniques to construct a rigorous argument. The solution must be self-contained and reference only the material from the course.

\vspace{0.5cm}

\end{document}